% Part: first-order-logic
% Chapter: proof-systems
% Section: seqeunt-calculus

\documentclass[../../../include/open-logic-section]{subfiles}

\begin{document}

\iftag{FOL}
      {\olfileid{fol}{prf}{seq}}
      {\olfileid{pl}{prf}{seq}}

\olsection{حساب التتابعيات}

ليست مادة حساب التتابعيات ترتيبًا من ال!!{sentence}s وحدها، كما في أنظمة
!!{derivation} كثيرة، بل \emph{تتابعيات}. والتتابعية على هذه الصورة:
\[
!A_\OLMathNumeral{1}, \dots, !A_\OLId{latin-ordinary}{m} \Sequent !B_\OLMathNumeral{1}, \dots, !B_\OLId{latin-ordinary}{n},
\]
فهي زوج متتاليتين من ال!!{sentence}s، بينهما الرمز~$\Sequent$؛ ولكل منهما أن تكون خالية.
\textbf{تنبيه تصحيحي على الأصل (C065-01).} جُعل مؤشر آخر صيغة في الطرف الأيمن
$\OLId{latin-ordinary}{n}$ مستقلًا عن المؤشر $\OLId{latin-ordinary}{m}$ في الطرف الأيسر؛ فليس يلزم أن يتساوى طولا
المتتاليتين.
و!!^a{derivation} في هذا الحساب شجرة تتابعيات، أعلاها تتابعيات مخصوصة تُسمى «ابتدائية» أو «بديهيات»، وما دونها ينتج كلٌّ منه عن التتابعيات التي تعلوه مباشرة بإحدى قواعد الاستدلال.
والقواعد على ضربين: ضرب يتصرف في ال!!{sentence}s بإضافة أو حذف أو تقديم وتأخير في أحد الجانبين؛ وضرب يُدخل
!!{formula} مركبة في النتيجة. من ذلك أن~$\LeftR{\land}$ تنقل من
$!A,
\OLId{greek-variable}{Gamma} \Sequent \OLId{greek-variable}{Delta}$ إلى $!A \land !B, \OLId{greek-variable}{Gamma} \Sequent \OLId{greek-variable}{Delta}$،
وأن~$\RightR{\lif}$ تنقل من $!A, \OLId{greek-variable}{Gamma} \Sequent
\OLId{greek-variable}{Delta}, !B$ إلى $\OLId{greek-variable}{Gamma} \Sequent \OLId{greek-variable}{Delta}, !A \lif !B$.
وهذا جائز لأي
$\OLId{greek-variable}{Gamma}$ و$\OLId{greek-variable}{Delta}$ و$!A$ و~$!B$، ولو كانت $\OLId{greek-variable}{Gamma}$
و~$\OLId{greek-variable}{Delta}$ خاليتين.

تُعرّف علاقة~$\Proves$ القائمة على حساب التتابعيات كما يأتي:
$\OLId{greek-variable}{Gamma} \Proves !A$ إذا وفقط إذا وجدت متتالية~$\OLId{greek-variable}{Gamma}_\OLMathNumeral{0}$ بحيث يكون
كل~$!A$ في~$\OLId{greek-variable}{Gamma}_\OLMathNumeral{0}$ عنصرًا في~$\OLId{greek-variable}{Gamma}$، ويوجد !!{derivation} تقع
التتابعية~$\OLId{greek-variable}{Gamma}_\OLMathNumeral{0} \Sequent !A$ في جذره. وتكون~$!A$ مبرهنة في حساب
التتابعيات إذا كان للتتابعية~$\Sequent !A$ !!a{derivation}. فمثلًا،
هذا !!a{derivation} يبين أن $\Proves (!A \land !B) \lif !A$:
\begin{prooftree}
  \Axiom$!A \fCenter !A$
  \RightLabel{\LeftR{\land}}
  \UnaryInf$!A \land !B \fCenter !A$
  \RightLabel{\RightR{\lif}}
  \UnaryInf$\fCenter (!A \land !B) \lif !A$
\end{prooftree}

أما عدم الاتساق فأن يوجد للمجموعة~$\OLId{greek-variable}{Gamma}$
!!a{derivation} لـ~$\OLId{greek-variable}{Gamma}_\OLMathNumeral{0} \Sequent$، حيث كل
$!A \in \OLId{greek-variable}{Gamma}_\OLMathNumeral{0}$ واقع في~$\OLId{greek-variable}{Gamma}$، والجانب الأيمن خالٍ.
وعندئذ تكفي قاعدة~\RightR{\Weakening} لجعل كل
!!{sentence} قابلة لل!!{derive} من تلك المجموعة.

وضع غيرهارد غنتسن حساب التتابعيات في ثلاثينيات القرن العشرين. وانتظام بنائه وتناظره يجعلان منه أداة صورية نافعة جدًا لنظرية ال!!{derivation}s. والعثور على
!!{derivation}s فيه يسير نسبيًا؛ غير أن قراءة هذه ال!!{derivation}s قد تعسر، وقد تخفى صلتها بالبراهين. ومع ذلك فقد اتضح حسن هذه الطريقة في أنظمة ال!!{derivation}، حتى صار لكثير من الأنساق المنطقية حساب تتابعيات خاص بها.

\end{document}
